\begin{pastexam}[label={px:df3-112-ds-04}]{地方三等資訊處理 112 年 第四題}
（申論；24 分）

針對如下的有向圖（節點為走訪對象，連線上的數字為走訪的 cost），依如下 BFS（配合 queue）與
DFS（配合 stack）演算法，進行所有節點的走訪，多個節點可以走訪時，以連線上 cost 較低者優先，
結果請以迴圈內部的顯示要求，依下表形式填入（stack 垂直表示，開口在上方，queue 水平表示，出口
在左，入口在右）。註：假設節點 S 為起始點。（24 分）

\begin{center}
\begin{tikzpicture}[x=1.05cm, y=1.05cm, >=Stealth,
  v/.style={circle, draw, minimum size=5.5mm, inner sep=0pt, font=\small},
  w/.style={font=\small, inner sep=1pt}]
\node[font=\small] at (0.2,0.35) {Directed Graph};
\node[v] (S) at (2,0) {S};
\node[v] (A) at (0.5,-1.3) {A};
\node[v] (B) at (2,-1.3) {B};
\node[v] (C) at (3.5,-1.3) {C};
\node[v] (X) at (0.5,-2.6) {X};
\node[v] (Y) at (2,-2.6) {Y};
\node[v] (Z) at (3.5,-2.6) {Z};
\draw[->] (S) -- node[w, above left, pos=0.5] {21} (A);
\draw[->] (S) -- node[w, left, pos=0.5] {11} (B);
\draw[->] (S) -- node[w, above right, pos=0.5] {31} (C);
\draw[->] (B) -- node[w, above] {1} (A);
\draw[->] (B) -- node[w, above] {4} (C);
\draw[->] (B) -- node[w, left, pos=0.4] {2} (X);
\draw[->] (B) -- node[w, left] {3} (Y);
\draw[->] (C) -- node[w, left] {5} (Z);
\draw[->] (Z) to[out=0, in=0, looseness=1.6] node[w, right] {10} (C);
\draw[->] (X) -- node[w, above] {6} (Y);
\draw[->] (Y) -- node[w, above] {8} (Z);
\draw[->] (X) to[out=-30, in=-150] node[w, fill=white] {7} (Z);
\draw[->] (Z.south) .. controls (3.1,-4.3) and (-0.9,-3.9) .. node[w, below, pos=0.6] {9} (A.south west);
\end{tikzpicture}
\end{center}
\begin{center}
\small
\begin{tabular}{|l|c|c|}
\hline
BFS 演算法 & Loop1 & $\cdots$ \\ \hline
print node & & \\ \hline
queue & \rule{0pt}{1.4cm} & \\ \hline
processSet & & \\ \hline
\end{tabular}
\quad
\begin{tabular}{|l|c|c|}
\hline
DFS 演算法 & Loop1 & $\cdots$ \\ \hline
print node & & \\ \hline
Stack & \rule{0pt}{1.4cm} & \\ \hline
processSet & & \\ \hline
\end{tabular}
\end{center}

BFS/DFS 演算法（/前為 BFS 使用 queue, /後為 DFS 使用 stack）
\begin{quote}
Step1: set \fbox{\underline{queue/stack}} to empty\\
\hspace*{3.2em}set \underline{processSet} to empty\\
Step2: \fbox{enqueue/push} \underline{S} and add \underline{S} into \underline{processSet}\\
Step3: while \fbox{\underline{queue/stack}} is not empty\\
Step31: \fbox{dequeue/pop} and print it\\
Step32: \fbox{enqueue/push} all one step neighbors which are not in \underline{processSet}\\
\hspace*{3.6em}according to the cost of edges and add them into \underline{processSet}\\
Step33: display content of \fbox{\underline{queue/stack}} and \underline{processSet}
\end{quote}
\end{pastexam}

\begin{pastsoln}
\textbf{先整理每個節點的出邊}，依 cost 由小到大：
\begin{center}
\small
\begin{tabular}{ll}
\toprule
節點 & 出邊（cost 小的在前） \\
\midrule
S & B (11), A (21), C (31) \\
B & A (1), X (2), Y (3), C (4) \\
C & Z (5) \\
X & Y (6), Z (7) \\
Y & Z (8) \\
Z & A (9), C (10) \\
A & （沒有出邊） \\
\bottomrule
\end{tabular}
\end{center}
這個演算法在\textbf{放進} queue／stack 時就把節點加入 processSet，所以每個節點只會被放進去一次。
「cost 較低者優先」對 queue 是依 cost 由小到大 enqueue（小的先出來）；對 stack 要讓 cost 小的
先被 pop，所以依 cost \textbf{由大到小} push，cost 最小的在最上面。 下面把題目的表轉成一列一個
Loop（欄位相同，只是轉置，方便排版）。 queue 左邊是出口；stack 照題目畫成直的，最上面是頂端
（開口）；processSet 依加入的先後寫。

\textbf{BFS（queue）}。 Step2 之後 queue $=$ S，processSet $= \{$S$\}$。
\begin{center}
\small
\begin{tabular}{cclp{4.2cm}}
\toprule
Loop & print node & queue（出口 $\leftarrow$ \quad $\leftarrow$ 入口） & processSet \\
\midrule
1 & S & B\ A\ C & S, B, A, C \\
2 & B & A\ C\ X\ Y & S, B, A, C, X, Y \\
3 & A & C\ X\ Y & S, B, A, C, X, Y \\
4 & C & X\ Y\ Z & S, B, A, C, X, Y, Z \\
5 & X & Y\ Z & S, B, A, C, X, Y, Z \\
6 & Y & Z & S, B, A, C, X, Y, Z \\
7 & Z & （空） & S, B, A, C, X, Y, Z \\
\bottomrule
\end{tabular}
\end{center}
Loop 2：B 的鄰居 A、C 已在 processSet，只 enqueue X (2)、Y (3)。 Loop 4：C 的鄰居 Z 不在，
enqueue。 Loop 5--7：X、Y、Z 的鄰居都已在 processSet。 \textbf{BFS 印出順序：S, B, A, C, X, Y, Z}。

\textbf{DFS（stack）}。 Step2 之後 stack $=$ S，processSet $= \{$S$\}$。
\begin{center}
\small
\renewcommand{\arraystretch}{1.15}
\begin{tabular}{cccp{4.2cm}}
\toprule
Loop & print node & stack（上面是頂端） & processSet \\
\midrule
1 & S & \begin{tabular}[c]{@{}c@{}}B\\A\\C\end{tabular} & S, C, A, B \\
2 & B & \begin{tabular}[c]{@{}c@{}}X\\Y\\A\\C\end{tabular} & S, C, A, B, Y, X \\
3 & X & \begin{tabular}[c]{@{}c@{}}Z\\Y\\A\\C\end{tabular} & S, C, A, B, Y, X, Z \\
4 & Z & \begin{tabular}[c]{@{}c@{}}Y\\A\\C\end{tabular} & S, C, A, B, Y, X, Z \\
5 & Y & \begin{tabular}[c]{@{}c@{}}A\\C\end{tabular} & S, C, A, B, Y, X, Z \\
6 & A & C & S, C, A, B, Y, X, Z \\
7 & C & （空） & S, C, A, B, Y, X, Z \\
\bottomrule
\end{tabular}
\end{center}
Loop 1：S 的鄰居依 cost 由大到小 push C (31)、A (21)、B (11)，B 在最上面。 Loop 2：pop B；
B 的鄰居 A、C 已在 processSet，依 cost 由大到小 push Y (3)、X (2)。 Loop 3：pop X；Y 已在，
push Z (7)。 Loop 4--7：Z、Y、A、C 的鄰居都已在 processSet，只有 pop。 \textbf{DFS 印出順序：
S, B, X, Z, Y, A, C}。

\textbf{另一種讀法}：若把 Step32 的「according to the cost」讀成 push 也依 cost 由小到大（cost
最大的在頂端），stack 依序是 S $\to$ (B A C) $\to$ (B A Z) $\to$ (B A) $\to$ (B) $\to$ (X Y)
$\to$ (X) $\to$ 空（頂端寫在右邊），印出 S, C, Z, A, B, Y, X；但這樣 cost 大的反而先走，和
題目「cost 較低者優先」不合，所以採用上面的寫法。

補充：題目的演算法在 push 時就標記，印出的順序不一定是真正的深度優先順序。 資結書第~21.4~節
的 DFS 是在\textbf{拜訪}時才標記，同樣從 S、cost 小的先走，順序是 S, B, A, X, Y, Z, C。 考試時
照題目給的步驟做即可。
\end{pastsoln}
